\chapter{Dynamic Programming: 1D DP}
\chapterepigraph{Dynamic programming is not a different way of thinking
from recursion -- it is recursion, plus a memory. Find the correct
recursive recurrence first, exactly as in the previous chapter; only
afterward ask whether that recurrence revisits the same sub-problem more
than once, and if it does, give it a memory.}

\section{From Recursion to Tabulation: The Four Stages}
\label{sec:dp-intro}

The Recursion chapter ended on its Word Break section, which exposed
a problem whose correct recursive solution recomputes
the same sub-problems exponentially many times. This chapter, and every
dynamic programming chapter that follows, is about systematically fixing
exactly that weakness. Every problem here is solved in up to four
progressively refined stages, and the chapter's single most important
habit is deriving them \emph{in this order}, never skipping straight to
the last one.

\begin{ideabox}[The Four Stages, Illustrated on Fibonacci]
Consider computing the $n$-th Fibonacci number, $F(n) = F(n-1) + F(n-2)$,
with $F(0)=0$, $F(1)=1$.

\begin{enumerate}
  \item \textbf{Recursion (brute force).} Translate the recurrence
        directly into a function: $\textsc{Fib}(n)$ returns $n$ if $n
        \le 1$, else $\textsc{Fib}(n-1) + \textsc{Fib}(n-2)$. This is
        always correct by definition, but $\textsc{Fib}(n-2)$ is
        computed once directly here and \emph{again} inside the
        recursive expansion of $\textsc{Fib}(n-1)$ -- and this
        duplication compounds at every level, giving $O(2^n)$ time.
  \item \textbf{Memoization (top-down DP).} Add an array (or hash map)
        $\textit{memo}$, indexed by the same parameter(s) the recursive
        function takes. Before computing $\textsc{Fib}(n)$'s result,
        check whether $\textit{memo}[n]$ is already filled; if so,
        return it directly. Otherwise compute it exactly as in the
        brute-force recursion, but \textbf{store} the result in
        $\textit{memo}[n]$ before returning. This is the \emph{exact
        same} recursive function, with two added lines -- yet it drops
        the time complexity to $O(n)$, since each distinct sub-problem
        $\textsc{Fib}(k)$ is now genuinely computed only once.
  \item \textbf{Tabulation (bottom-up DP).} Rewrite the same recurrence
        as an explicit loop filling an array $\textit{dp}$ from the base
        cases \emph{upward} to $n$, instead of recursing downward from
        $n$: $\textit{dp}[0]=0$, $\textit{dp}[1]=1$, and for $i$ from $2$
        to $n$: $\textit{dp}[i] = \textit{dp}[i-1] + \textit{dp}[i-2]$.
        This computes exactly the same values as memoization, in the
        same $O(n)$ time, but with no recursion or call-stack overhead
        at all.
  \item \textbf{Space Optimization.} Observe that computing
        $\textit{dp}[i]$ only ever needs $\textit{dp}[i-1]$ and
        $\textit{dp}[i-2]$ -- once $\textit{dp}[i]$ is computed, every
        earlier entry is permanently unneeded. Replace the full array
        with two scalar variables, updated in a rolling fashion, cutting
        space from $O(n)$ to $O(1)$.
\end{enumerate}
\end{ideabox}

\subsection*{Why Derive All Four, Rather Than Jumping to the Last}

Tabulation and space optimization are mechanical rewrites of memoization,
which is itself a two-line addition to brute-force recursion -- but each
rewrite is only safe once the \emph{previous} stage is verified correct.
Attempting to write a tabulation loop directly, without first confirming
the recursive recurrence by hand on a small example, is the single most
common source of off-by-one and base-case bugs in dynamic programming.
This chapter therefore presents every problem in exactly this order:
state the recursive recurrence and verify it on a small case, add
memoization, convert to tabulation, then space-optimize -- and the
\emph{only} thing that changes from problem to problem is what the
recurrence itself looks like.

\begin{patternbox}[How to Recognise a 1D DP Problem and Find Its Recurrence]
\begin{itemize}
  \item \textbf{``Number of ways to reach step $n$,'' or ``minimum
        cost/jumps to reach position $n$''} -- the answer at position
        $n$ is built from the answers at a small, fixed set of earlier
        positions (e.g.\ $n-1$ and $n-2$ for single or double steps),
        combined by addition (counting ways) or by minimization (cost).
  \item \textbf{``Maximum sum, no two adjacent chosen'' (House Robber and
        its variants)} -- at each position, the recurrence is a
        \emph{choice}: include this position's value plus the best
        answer from two positions back, or exclude it and take the best
        answer from one position back, whichever is larger.
\end{itemize}
Once the recurrence is identified, the parameter(s) it depends on become
the memoization array's index (or indices) -- this is true for every DP
problem in this book, not just the ones in this chapter.
\end{patternbox}

Each section in this chapter follows this four-stage structure, together
with the keyword/pattern recognition, a hand-traced dry run of the
tabulation table being filled, and a time/space complexity comparison
across all four stages.

\newpage

\section{Fibonacci Number}
\label{sec:dp-fibonacci}

\problemstatement{Compute the $n$-th Fibonacci number, defined by
$F(0)=0$, $F(1)=1$, and $F(n) = F(n-1) + F(n-2)$ for $n \ge 2$.}

\begin{patternbox}
\emph{``Each term is defined by combining a small, fixed number of
earlier terms''} is the simplest possible 1D DP keyword pattern, and
Fibonacci is its archetype: the recurrence $F(n) = F(n-1)+F(n-2)$ is
already explicitly given, so no derivation is needed -- only the
discipline of carrying it through all four stages described in
Section~\ref{sec:dp-intro}.
\end{patternbox}

\subsection*{Approach 1: Recursion}

The recurrence translates directly into a function with two base cases
($n=0$ and $n=1$) and one recursive case.

\begin{lstlisting}
#include <bits/stdc++.h>
using namespace std;

int fibRecursive(int n) {
    if (n <= 1) return n;
    return fibRecursive(n - 1) + fibRecursive(n - 2);
}

int main() {
    cout << fibRecursive(10) << endl; // 55
    return 0;
}
\end{lstlisting}

\begin{complexitybox}[Approach 1 Complexity]
\textbf{Time.} Each call spawns two more, and the recursion tree has
depth $n$, giving $O(2^n)$ -- more precisely $O(\phi^n)$ where $\phi$ is
the golden ratio, but exponential either way.
\textbf{Space.} $O(n)$ for the recursion stack depth.
\end{complexitybox}

\subsection*{Approach 2: Memoization (Top-Down DP)}

\begin{ideabox}[Why $\textsc{Fib}(k)$ Is Computed Repeatedly]
Expanding $\textsc{Fib}(5) = \textsc{Fib}(4) + \textsc{Fib}(3)$: the call
to $\textsc{Fib}(4)$ itself expands to $\textsc{Fib}(3)+\textsc{Fib}(2)$,
so $\textsc{Fib}(3)$ is computed once directly and once more inside
$\textsc{Fib}(4)$'s expansion -- and this doubling compounds at every
level. Caching each $\textsc{Fib}(k)$'s result the first time it is
computed, in an array $\textit{memo}$ indexed by $k$, means every later
request for the same $k$ is an $O(1)$ lookup instead of a full
re-expansion.
\end{ideabox}

\begin{lstlisting}
#include <bits/stdc++.h>
using namespace std;

int fibMemo(int n, vector<int>& memo) {
    if (n <= 1) return n;
    if (memo[n] != -1) return memo[n];

    memo[n] = fibMemo(n - 1, memo) + fibMemo(n - 2, memo);
    return memo[n];
}

int main() {
    int n = 10;
    vector<int> memo(n + 1, -1);
    cout << fibMemo(n, memo) << endl; // 55
    return 0;
}
\end{lstlisting}

\begin{complexitybox}[Approach 2 Complexity]
\textbf{Time.} Each of the $n+1$ distinct sub-problems
$\textsc{Fib}(0),\ldots,\textsc{Fib}(n)$ is computed exactly once,
each doing $O(1)$ work beyond its (now-cached) recursive calls, giving
$O(n)$.
\textbf{Space.} $O(n)$ for the memo array, plus $O(n)$ for the
recursion stack -- still $O(n)$ overall, but for a different reason than
Approach 1 (here it's the cache, not the redundant branching, that
dominates the asymptotic class, even though the stack depth alone was
already $O(n)$ in Approach 1 too).
\end{complexitybox}

\subsection*{Approach 3: Tabulation (Bottom-Up DP)}

\begin{ideabox}[Filling the Table in the Opposite Direction]
Memoization computes $\textsc{Fib}(n)$ by recursing \emph{downward}
toward the base cases, caching as it unwinds. Tabulation instead starts
\emph{at} the base cases and fills an array $\textit{dp}$ \emph{upward}
toward $n$ in an explicit loop -- computing the exact same values, in the
same dependency order, but without any function call overhead.
\end{ideabox}

\begin{lstlisting}
#include <bits/stdc++.h>
using namespace std;

int fibTabulation(int n) {
    if (n == 0) return 0;

    vector<int> dp(n + 1);
    dp[0] = 0;
    dp[1] = 1;
    for (int i = 2; i <= n; i++) {
        dp[i] = dp[i - 1] + dp[i - 2];
    }
    return dp[n];
}

int main() {
    cout << fibTabulation(10) << endl; // 55
    return 0;
}
\end{lstlisting}

\subsection*{Dry run of the tabulation table}

\dryrun{Take $n=6$.}

\begin{center}
\begin{tabular}{c|c|c|c|c|c|c|c}
$i$ & 0 & 1 & 2 & 3 & 4 & 5 & 6 \\ \hline
$\textit{dp}[i]$ & 0 & 1 & 1 & 2 & 3 & 5 & 8
\end{tabular}
\end{center}

Each entry from $i=2$ onward is filled as $\textit{dp}[i-1]+\textit{dp}[i-2]$:
e.g.\ $\textit{dp}[4] = \textit{dp}[3]+\textit{dp}[2] = 2+1 = 3$. The
answer, $\textit{dp}[6]=8$, matches $F(6)=8$.

\begin{complexitybox}[Approach 3 Complexity]
\textbf{Time.} $O(n)$ -- one pass filling $n-1$ table entries.
\textbf{Space.} $O(n)$ for the $\textit{dp}$ array; \textbf{no}
recursion stack at all, unlike Approaches 1 and 2.
\end{complexitybox}

\subsection*{Approach 4: Space Optimization}

\begin{ideabox}[Only the Last Two Entries Are Ever Needed]
Computing $\textit{dp}[i]$ only reads $\textit{dp}[i-1]$ and
$\textit{dp}[i-2]$ -- every entry before $i-2$ is permanently dead the
moment $\textit{dp}[i]$ is computed. Replacing the full array with two
rolling scalars, updated each iteration, preserves correctness while
cutting space from $O(n)$ to $O(1)$.
\end{ideabox}

\begin{lstlisting}
#include <bits/stdc++.h>
using namespace std;

int fibSpaceOptimized(int n) {
    if (n == 0) return 0;

    int prev2 = 0, prev1 = 1;
    for (int i = 2; i <= n; i++) {
        int current = prev1 + prev2;
        prev2 = prev1;
        prev1 = current;
    }
    return prev1;
}

int main() {
    cout << fibSpaceOptimized(10) << endl; // 55
    return 0;
}
\end{lstlisting}

\begin{complexitybox}[Approach 4 Complexity]
\textbf{Time.} $O(n)$, unchanged from tabulation.
\textbf{Space.} $O(1)$ -- only two scalar variables, regardless of how
large $n$ is.
\end{complexitybox}

\begin{takeawaybox}
Every later 1D DP problem in this chapter follows exactly this
progression. The only thing that will change from here on is the
\emph{recurrence itself} -- how many earlier states each step depends on,
and how they are combined (summed, for counting; minimized or maximized,
for optimization). The four-stage scaffolding shown here is reused
verbatim.
\end{takeawaybox}

\section{Climbing Stairs}
\label{sec:dp-climbing-stairs}

\problemstatement{A staircase has $n$ steps. From any step, you may
climb either $1$ or $2$ steps at a time. Count the number of distinct
ways to reach the top (step $n$), starting from step $0$.}

\begin{patternbox}
\emph{``Number of distinct ways to reach step $n$, taking $1$ or $2$
steps at a time''} is the direct counting analogue of
Section~\ref{sec:dp-fibonacci}'s recurrence: the last move to reach step
$n$ was either a single step from $n-1$, or a double step from $n-2$ --
so every way of reaching $n$ is exactly one way of reaching $n-1$,
extended by one step, \emph{or} one way of reaching $n-2$, extended by
two steps, and these two groups are disjoint and exhaustive.
\end{patternbox}

\subsection*{Deriving the recurrence}

Let $\textit{ways}(i)$ be the number of distinct ways to reach step $i$.
Since the only moves are $+1$ and $+2$, the last move into step $i$ came
from step $i-1$ or step $i-2$ -- so $\textit{ways}(i) =
\textit{ways}(i-1) + \textit{ways}(i-2)$. Base cases:
$\textit{ways}(0)=1$ (one way to be at the start: do nothing), and
$\textit{ways}(1)=1$ (only one way to reach step $1$: a single step from
$0$). This is exactly Fibonacci's recurrence, shifted by one index, with
different base cases.

\subsection*{Approach 1: Recursion}

\begin{lstlisting}
#include <bits/stdc++.h>
using namespace std;

int climbStairsRecursive(int n) {
    if (n <= 1) return 1;
    return climbStairsRecursive(n - 1) + climbStairsRecursive(n - 2);
}

int main() {
    cout << climbStairsRecursive(5) << endl; // 8
    return 0;
}
\end{lstlisting}

\begin{complexitybox}[Approach 1 Complexity]
\textbf{Time.} $O(2^n)$ -- identical overlapping-subproblem blowup as
plain Fibonacci.
\textbf{Space.} $O(n)$ for the recursion stack.
\end{complexitybox}

\subsection*{Approach 2: Memoization}

\begin{lstlisting}
#include <bits/stdc++.h>
using namespace std;

int climbStairsMemo(int n, vector<int>& memo) {
    if (n <= 1) return 1;
    if (memo[n] != -1) return memo[n];

    memo[n] = climbStairsMemo(n - 1, memo) + climbStairsMemo(n - 2, memo);
    return memo[n];
}

int main() {
    int n = 5;
    vector<int> memo(n + 1, -1);
    cout << climbStairsMemo(n, memo) << endl; // 8
    return 0;
}
\end{lstlisting}

\begin{complexitybox}[Approach 2 Complexity]
\textbf{Time.} $O(n)$.
\textbf{Space.} $O(n)$ memo array $+$ $O(n)$ recursion stack.
\end{complexitybox}

\subsection*{Approach 3: Tabulation}

\begin{lstlisting}
#include <bits/stdc++.h>
using namespace std;

int climbStairsTabulation(int n) {
    if (n <= 1) return 1;

    vector<int> dp(n + 1);
    dp[0] = 1;
    dp[1] = 1;
    for (int i = 2; i <= n; i++) {
        dp[i] = dp[i - 1] + dp[i - 2];
    }
    return dp[n];
}

int main() {
    cout << climbStairsTabulation(5) << endl; // 8
    return 0;
}
\end{lstlisting}

\subsection*{Dry run of the tabulation table}

\dryrun{Take $n=5$.}

\begin{center}
\begin{tabular}{c|c|c|c|c|c|c}
$i$ & 0 & 1 & 2 & 3 & 4 & 5 \\ \hline
$\textit{dp}[i]$ & 1 & 1 & 2 & 3 & 5 & 8
\end{tabular}
\end{center}

$\textit{dp}[5]=8$: indeed, there are $8$ distinct ways to climb $5$
steps using moves of size $1$ or $2$ (e.g.\ $1{+}1{+}1{+}1{+}1$,
$2{+}1{+}1{+}1$ in every position, $2{+}2{+}1$ in every order, and so
on).

\begin{complexitybox}[Approach 3 Complexity]
\textbf{Time.} $O(n)$.
\textbf{Space.} $O(n)$ for the table, no recursion stack.
\end{complexitybox}

\subsection*{Approach 4: Space Optimization}

\begin{lstlisting}
#include <bits/stdc++.h>
using namespace std;

int climbStairsSpaceOptimized(int n) {
    if (n <= 1) return 1;

    int prev2 = 1, prev1 = 1;
    for (int i = 2; i <= n; i++) {
        int current = prev1 + prev2;
        prev2 = prev1;
        prev1 = current;
    }
    return prev1;
}

int main() {
    cout << climbStairsSpaceOptimized(5) << endl; // 8
    return 0;
}
\end{lstlisting}

\begin{complexitybox}[Approach 4 Complexity]
\textbf{Time.} $O(n)$.
\textbf{Space.} $O(1)$.
\end{complexitybox}

\begin{takeawaybox}
Recognizing that ``count the ways to reach $n$ via moves from a small
fixed set'' always produces a recurrence that \textbf{sums} the way-counts
of the positions those moves originate from is the key insight -- the
specific move set (here, $\{1,2\}$) only determines \emph{how many}
earlier terms the recurrence sums, and which ones. Section~\ref{sec:dp-frog-jump}
generalizes this from counting ways to minimizing cost, with the sum
replaced by a minimum.
\end{takeawaybox}

\section{Frog Jump}
\label{sec:dp-frog-jump}

\problemstatement{A frog starts at index $0$ of an array $\textit{height}$
and wants to reach index $n-1$. From index $i$, it may jump to $i+1$ or
$i+2$, at a cost of $|\textit{height}[i] - \textit{height}[j]|$ for a
jump to index $j$. Find the minimum total cost to reach index $n-1$.}

\begin{patternbox}
\emph{``Minimum cost to reach position $n$, moving via a small fixed set
of jump sizes''} is Section~\ref{sec:dp-climbing-stairs}'s recurrence
shape with \textbf{sum replaced by minimum}: instead of \emph{counting}
every way to combine smaller answers, we \emph{minimize} over them,
because here we want the cheapest path, not a tally of all paths.
\end{patternbox}

\subsection*{Deriving the recurrence}

Let $\textit{minCost}(i)$ be the minimum cost to reach index $i$ from
index $0$. The frog's last jump into index $i$ came from either $i-1$
(cost $|\textit{height}[i]-\textit{height}[i-1]|$ added to
$\textit{minCost}(i-1)$) or $i-2$ (cost
$|\textit{height}[i]-\textit{height}[i-2]|$ added to
$\textit{minCost}(i-2)$) -- and since we want the \emph{cheapest} way to
reach $i$, we take the minimum of these two options:
\[
  \textit{minCost}(i) = \min\big(\textit{minCost}(i-1) +
  |\textit{height}[i]-\textit{height}[i-1]|,\ \textit{minCost}(i-2) +
  |\textit{height}[i]-\textit{height}[i-2]|\big).
\]
Base case: $\textit{minCost}(0)=0$ (no cost to start where you already
are); for $i=1$, only the jump from $0$ is possible (there is no index
$-1$), so $\textit{minCost}(1) = |\textit{height}[1]-\textit{height}[0]|$
directly.

\subsection*{Approach 1: Recursion}

\begin{lstlisting}
#include <bits/stdc++.h>
using namespace std;

int frogJumpRecursive(int i, vector<int>& height) {
    if (i == 0) return 0;
    if (i == 1) return abs(height[1] - height[0]);

    int fromOneBack = frogJumpRecursive(i - 1, height) + abs(height[i] - height[i-1]);
    int fromTwoBack  = frogJumpRecursive(i - 2, height) + abs(height[i] - height[i-2]);
    return min(fromOneBack, fromTwoBack);
}

int main() {
    vector<int> height = {10, 30, 40, 50};
    int n = height.size();
    cout << frogJumpRecursive(n - 1, height) << endl; // 40
    return 0;
}
\end{lstlisting}

\begin{complexitybox}[Approach 1 Complexity]
\textbf{Time.} $O(2^n)$ -- the same two-branches-per-call structure as
every recurrence in this chapter so far.
\textbf{Space.} $O(n)$ recursion stack.
\end{complexitybox}

\subsection*{Approach 2: Memoization}

\begin{lstlisting}
#include <bits/stdc++.h>
using namespace std;

int frogJumpMemo(int i, vector<int>& height, vector<int>& memo) {
    if (i == 0) return 0;
    if (i == 1) return abs(height[1] - height[0]);
    if (memo[i] != -1) return memo[i];

    int fromOneBack = frogJumpMemo(i - 1, height, memo) + abs(height[i] - height[i-1]);
    int fromTwoBack  = frogJumpMemo(i - 2, height, memo) + abs(height[i] - height[i-2]);
    memo[i] = min(fromOneBack, fromTwoBack);
    return memo[i];
}

int main() {
    vector<int> height = {10, 30, 40, 50};
    int n = height.size();
    vector<int> memo(n, -1);
    cout << frogJumpMemo(n - 1, height, memo) << endl; // 40
    return 0;
}
\end{lstlisting}

\begin{complexitybox}[Approach 2 Complexity]
\textbf{Time.} $O(n)$.
\textbf{Space.} $O(n)$ memo $+$ $O(n)$ recursion stack.
\end{complexitybox}

\subsection*{Approach 3: Tabulation}

\begin{lstlisting}
#include <bits/stdc++.h>
using namespace std;

int frogJumpTabulation(vector<int>& height) {
    int n = height.size();
    vector<int> dp(n);
    dp[0] = 0;
    if (n > 1) dp[1] = abs(height[1] - height[0]);

    for (int i = 2; i < n; i++) {
        int fromOneBack = dp[i - 1] + abs(height[i] - height[i-1]);
        int fromTwoBack  = dp[i - 2] + abs(height[i] - height[i-2]);
        dp[i] = min(fromOneBack, fromTwoBack);
    }
    return dp[n - 1];
}

int main() {
    vector<int> height = {10, 30, 40, 50};
    cout << frogJumpTabulation(height) << endl; // 40
    return 0;
}
\end{lstlisting}

\subsection*{Dry run of the tabulation table}

\dryrun{Take $\textit{height} = [10,30,40,50]$.}

\begin{itemize}
  \item $\textit{dp}[0] = 0$.
  \item $\textit{dp}[1] = |30-10| = 20$.
  \item $\textit{dp}[2] = \min(\textit{dp}[1]+|40{-}30|,\
        \textit{dp}[0]+|40{-}10|) = \min(20{+}10,\ 0{+}30) =
        \min(30,30) = 30$.
  \item $\textit{dp}[3] = \min(\textit{dp}[2]+|50{-}40|,\
        \textit{dp}[1]+|50{-}30|) = \min(30{+}10,\ 20{+}20) =
        \min(40,40) = 40$.
\end{itemize}

\begin{center}
\begin{tabular}{c|c|c|c|c}
$i$ & 0 & 1 & 2 & 3 \\ \hline
$\textit{height}[i]$ & 10 & 30 & 40 & 50 \\ \hline
$\textit{dp}[i]$ & 0 & 20 & 30 & 40
\end{tabular}
\end{center}

The answer is $\textit{dp}[3]=40$: one cheapest path is
$0\to1\to2\to3$, costing $20+10+10=40$ (and the alternative
$0\to2\to3$, costing $30+10=40$, ties exactly).

\begin{complexitybox}[Approach 3 Complexity]
\textbf{Time.} $O(n)$.
\textbf{Space.} $O(n)$ table, no recursion stack.
\end{complexitybox}

\subsection*{Approach 4: Space Optimization}

\begin{lstlisting}
#include <bits/stdc++.h>
using namespace std;

int frogJumpSpaceOptimized(vector<int>& height) {
    int n = height.size();
    if (n == 1) return 0;

    int prev2 = 0, prev1 = abs(height[1] - height[0]);
    for (int i = 2; i < n; i++) {
        int fromOneBack = prev1 + abs(height[i] - height[i-1]);
        int fromTwoBack  = prev2 + abs(height[i] - height[i-2]);
        int current = min(fromOneBack, fromTwoBack);
        prev2 = prev1;
        prev1 = current;
    }
    return prev1;
}

int main() {
    vector<int> height = {10, 30, 40, 50};
    cout << frogJumpSpaceOptimized(height) << endl; // 40
    return 0;
}
\end{lstlisting}

\begin{complexitybox}[Approach 4 Complexity]
\textbf{Time.} $O(n)$.
\textbf{Space.} $O(1)$.
\end{complexitybox}

\begin{takeawaybox}
Swapping a sum for a minimum (or maximum) in the combination step is
often the \emph{only} change needed to turn a ``counting'' DP recurrence
into a ``cost optimization'' one, with every other part of the four-stage
derivation -- base cases, indexing, the eventual space optimization --
carrying over in the same shape. Section~\ref{sec:dp-frog-jump-k}
generalizes the jump set from $\{1,2\}$ to $\{1,\ldots,k\}$, widening the
minimum from two candidates to $k$.
\end{takeawaybox}

\section{Frog Jump with K Distances}
\label{sec:dp-frog-jump-k}

\problemstatement{Same setting as Section~\ref{sec:dp-frog-jump}, except
the frog may jump from index $i$ to any of $i+1, i+2, \ldots, i+k$ (up
to $k$ steps ahead, instead of only $1$ or $2$), at cost
$|\textit{height}[i]-\textit{height}[j]|$ for a jump to index $j$. Find
the minimum total cost to reach index $n-1$.}

\begin{patternbox}
\emph{``Same recurrence shape, but the fixed set of options widens from
$2$ to a parameter $k$''} -- recognizing this is the entire content of
this section: Section~\ref{sec:dp-frog-jump}'s minimum over $2$ candidates
($i-1$ and $i-2$) becomes a minimum over up to $k$ candidates
($i-1, i-2, \ldots, i-k$), with everything else about the derivation
unchanged.
\end{patternbox}

\subsection*{Deriving the recurrence}

The frog's last jump into index $i$ came from some index $i-j$ for some
$j$ between $1$ and $k$ (as long as $i-j \ge 0$). So:
\[
  \textit{minCost}(i) = \min_{1 \le j \le k,\ i-j \ge 0}
  \Big(\textit{minCost}(i-j) + |\textit{height}[i]-\textit{height}[i-j]|\Big).
\]
Base case: $\textit{minCost}(0)=0$.

\subsection*{Approach 1: Recursion}

\begin{lstlisting}
#include <bits/stdc++.h>
using namespace std;

int frogJumpKRecursive(int i, int k, vector<int>& height) {
    if (i == 0) return 0;

    int best = INT_MAX;
    for (int j = 1; j <= k && i - j >= 0; j++) {
        int cost = frogJumpKRecursive(i - j, k, height) +
                    abs(height[i] - height[i - j]);
        best = min(best, cost);
    }
    return best;
}

int main() {
    vector<int> height = {10, 30, 40, 50, 20};
    int n = height.size(), k = 3;
    cout << frogJumpKRecursive(n - 1, k, height) << endl;
    return 0;
}
\end{lstlisting}

\begin{complexitybox}[Approach 1 Complexity]
\textbf{Time.} $O(k^n)$ in the worst case -- each call now branches
into up to $k$ recursive calls instead of $2$.
\textbf{Space.} $O(n)$ for the recursion stack depth.
\end{complexitybox}

\subsection*{Approach 2: Memoization}

\begin{lstlisting}
#include <bits/stdc++.h>
using namespace std;

int frogJumpKMemo(int i, int k, vector<int>& height, vector<int>& memo) {
    if (i == 0) return 0;
    if (memo[i] != -1) return memo[i];

    int best = INT_MAX;
    for (int j = 1; j <= k && i - j >= 0; j++) {
        int cost = frogJumpKMemo(i - j, k, height, memo) +
                    abs(height[i] - height[i - j]);
        best = min(best, cost);
    }
    return memo[i] = best;
}

int main() {
    vector<int> height = {10, 30, 40, 50, 20};
    int n = height.size(), k = 3;
    vector<int> memo(n, -1);
    cout << frogJumpKMemo(n - 1, k, height, memo) << endl;
    return 0;
}
\end{lstlisting}

\begin{complexitybox}[Approach 2 Complexity]
\textbf{Time.} $O(n \cdot k)$ -- $n$ distinct sub-problems, each doing
$O(k)$ work in its loop.
\textbf{Space.} $O(n)$ memo $+$ $O(n)$ recursion stack.
\end{complexitybox}

\subsection*{Approach 3: Tabulation}

\begin{lstlisting}
#include <bits/stdc++.h>
using namespace std;

int frogJumpKTabulation(vector<int>& height, int k) {
    int n = height.size();
    vector<int> dp(n, 0);

    for (int i = 1; i < n; i++) {
        int best = INT_MAX;
        for (int j = 1; j <= k && i - j >= 0; j++) {
            int cost = dp[i - j] + abs(height[i] - height[i - j]);
            best = min(best, cost);
        }
        dp[i] = best;
    }
    return dp[n - 1];
}

int main() {
    vector<int> height = {10, 30, 40, 50, 20};
    cout << frogJumpKTabulation(height, 3) << endl;
    return 0;
}
\end{lstlisting}

\subsection*{Dry run of the tabulation table}

\dryrun{Take $\textit{height} = [10,30,40,50,20]$, $k=3$.}

\begin{itemize}
  \item $\textit{dp}[0]=0$.
  \item $\textit{dp}[1] = \textit{dp}[0]+|30{-}10| = 20$ (only $j=1$
        possible).
  \item $\textit{dp}[2] = \min(\textit{dp}[1]+|40{-}30|,\
        \textit{dp}[0]+|40{-}10|) = \min(30,30) = 30$.
  \item $\textit{dp}[3] = \min(\textit{dp}[2]+|50{-}40|,\
        \textit{dp}[1]+|50{-}30|,\ \textit{dp}[0]+|50{-}10|) =
        \min(40,40,40) = 40$.
  \item $\textit{dp}[4] = \min(\textit{dp}[3]+|20{-}50|,\
        \textit{dp}[2]+|20{-}40|,\ \textit{dp}[1]+|20{-}30|) =
        \min(40{+}30,\ 30{+}20,\ 20{+}10) = \min(70,50,30) = 30$.
\end{itemize}

\begin{center}
\begin{tabular}{c|c|c|c|c|c}
$i$ & 0 & 1 & 2 & 3 & 4 \\ \hline
$\textit{dp}[i]$ & 0 & 20 & 30 & 40 & 30
\end{tabular}
\end{center}

The answer is $\textit{dp}[4]=30$, achieved by jumping directly from
index $1$ ($\textit{height}=30$) to index $4$ ($\textit{height}=20$), a
jump of size $3 \le k$, costing $\textit{dp}[1] + |20-30| = 20+10=30$.

\begin{complexitybox}[Approach 3 Complexity]
\textbf{Time.} $O(n \cdot k)$.
\textbf{Space.} $O(n)$ table, no recursion stack.
\end{complexitybox}

\subsection*{Approach 4: Space Optimization}

\begin{ideabox}[Why This One Cannot Be Reduced All the Way to $O(1)$]
Unlike Section~\ref{sec:dp-frog-jump}, where only the \emph{two}
immediately preceding entries were ever needed, here computing
$\textit{dp}[i]$ may need any of the previous $k$ entries -- so the
space optimization here reduces storage to a \textbf{sliding window} of
the last $k$ entries (e.g.\ a small circular buffer or deque of size
$k$), rather than collapsing all the way to two scalar variables. When
$k$ is a fixed small constant, this is still $O(k) = O(1)$ relative to
$n$; when $k$ can grow proportionally with $n$, no reduction below
$O(n)$ is possible, since up to $n$ of the entries might genuinely be
needed.
\end{ideabox}

\begin{lstlisting}
#include <bits/stdc++.h>
using namespace std;

int frogJumpKSpaceOptimized(vector<int>& height, int k) {
    int n = height.size();
    vector<int> window(k, 0); // window[i % k] holds dp[i]
    window[0] = 0;

    for (int i = 1; i < n; i++) {
        int best = INT_MAX;
        for (int j = 1; j <= k && i - j >= 0; j++) {
            int cost = window[(i - j) % k] + abs(height[i] - height[i - j]);
            best = min(best, cost);
        }
        window[i % k] = best;
    }
    return window[(n - 1) % k];
}

int main() {
    vector<int> height = {10, 30, 40, 50, 20};
    cout << frogJumpKSpaceOptimized(height, 3) << endl;
    return 0;
}
\end{lstlisting}

\begin{complexitybox}[Approach 4 Complexity]
\textbf{Time.} $O(n \cdot k)$, unchanged.
\textbf{Space.} $O(k)$ for the sliding window, instead of $O(n)$ for
the full table.
\end{complexitybox}

\begin{takeawaybox}
Generalizing a recurrence's fixed-size option set from a small constant
(here, $2$, in Frog Jump) to a parameter $k$ widens every stage's loop
from $2$ iterations to $k$, multiplying the time complexity by $k$ and
capping how far the space optimization can shrink: only the last $k$
table entries, not just the last $2$, must remain available, so the
``two rolling variables'' trick of earlier sections becomes a
``sliding window of size $k$'' trick here.
\end{takeawaybox}

\section{House Robber}
\label{sec:dp-house-robber}

\problemstatement{Given an array $\textit{nums}$ where
$\textit{nums}[i]$ is the amount of money in house $i$, find the maximum
total amount that can be robbed without ever robbing two
\textbf{adjacent} houses.}

\begin{patternbox}
\emph{``Maximum sum, no two adjacent elements chosen''} is a different
recurrence \emph{shape} from every problem so far in this chapter: those
were all ``combine answers from a fixed set of earlier positions''
(summed, for counting, or minimized, for cost). This one is a genuine
\textbf{include-or-exclude choice} at every position -- structurally the
same choice as the pick/not-pick recursion from
the Recursion chapter's Power Set section, but now resolved by \textbf{maximizing}
rather than branching into both options forever.
\end{patternbox}

\subsection*{Deriving the recurrence}

Let $f(i)$ be the maximum amount robbable considering only houses $0$
through $i$. At house $i$, there are exactly two choices: \textbf{rob}
house $i$ (gaining $\textit{nums}[i]$, but then house $i-1$ \emph{must}
be skipped, so we add the best answer from $f(i-2)$), or \textbf{skip}
house $i$ entirely (carrying forward $f(i-1)$ unchanged, since skipping
loses nothing already secured). We take whichever choice is larger:
\[
  f(i) = \max\big(f(i-1),\ f(i-2) + \textit{nums}[i]\big).
\]
Base cases: $f(0) = \textit{nums}[0]$ (only one house, must take it for
the best possible answer at this point), and, treating an empty prefix
before index $0$ as contributing $0$, we can write $f(-1)=0$ for
notational convenience when computing $f(1) = \max(f(0), f(-1) +
\textit{nums}[1]) = \max(\textit{nums}[0], \textit{nums}[1])$.

\subsection*{Approach 1: Recursion}

\begin{lstlisting}
#include <bits/stdc++.h>
using namespace std;

int robRecursive(int i, vector<int>& nums) {
    if (i == 0) return nums[0];
    if (i < 0) return 0;

    int skip = robRecursive(i - 1, nums);
    int take = robRecursive(i - 2, nums) + nums[i];
    return max(skip, take);
}

int main() {
    vector<int> nums = {1, 2, 3, 1};
    int n = nums.size();
    cout << robRecursive(n - 1, nums) << endl; // 4
    return 0;
}
\end{lstlisting}

\begin{complexitybox}[Approach 1 Complexity]
\textbf{Time.} $O(2^n)$.
\textbf{Space.} $O(n)$ recursion stack.
\end{complexitybox}

\subsection*{Approach 2: Memoization}

\begin{lstlisting}
#include <bits/stdc++.h>
using namespace std;

int robMemo(int i, vector<int>& nums, vector<int>& memo) {
    if (i == 0) return nums[0];
    if (i < 0) return 0;
    if (memo[i] != -1) return memo[i];

    int skip = robMemo(i - 1, nums, memo);
    int take = robMemo(i - 2, nums, memo) + nums[i];
    return memo[i] = max(skip, take);
}

int main() {
    vector<int> nums = {1, 2, 3, 1};
    int n = nums.size();
    vector<int> memo(n, -1);
    cout << robMemo(n - 1, nums, memo) << endl; // 4
    return 0;
}
\end{lstlisting}

\begin{complexitybox}[Approach 2 Complexity]
\textbf{Time.} $O(n)$.
\textbf{Space.} $O(n)$ memo $+$ $O(n)$ recursion stack.
\end{complexitybox}

\subsection*{Approach 3: Tabulation}

\begin{lstlisting}
#include <bits/stdc++.h>
using namespace std;

int robTabulation(vector<int>& nums) {
    int n = nums.size();
    vector<int> dp(n);
    dp[0] = nums[0];

    for (int i = 1; i < n; i++) {
        int skip = dp[i - 1];
        int take = (i >= 2 ? dp[i - 2] : 0) + nums[i];
        dp[i] = max(skip, take);
    }
    return dp[n - 1];
}

int main() {
    vector<int> nums = {1, 2, 3, 1};
    cout << robTabulation(nums) << endl; // 4
    return 0;
}
\end{lstlisting}

\subsection*{Dry run of the tabulation table}

\dryrun{Take $\textit{nums} = [1,2,3,1]$.}

\begin{itemize}
  \item $\textit{dp}[0] = 1$.
  \item $\textit{dp}[1] = \max(\textit{dp}[0],\ 0+2) = \max(1,2) = 2$.
  \item $\textit{dp}[2] = \max(\textit{dp}[1],\ \textit{dp}[0]+3) =
        \max(2,\ 1+3) = \max(2,4) = 4$.
  \item $\textit{dp}[3] = \max(\textit{dp}[2],\ \textit{dp}[1]+1) =
        \max(4,\ 2+1) = \max(4,3) = 4$.
\end{itemize}

\begin{center}
\begin{tabular}{c|c|c|c|c}
$i$ & 0 & 1 & 2 & 3 \\ \hline
$\textit{nums}[i]$ & 1 & 2 & 3 & 1 \\ \hline
$\textit{dp}[i]$ & 1 & 2 & 4 & 4
\end{tabular}
\end{center}

The answer is $\textit{dp}[3]=4$, achieved by robbing houses $0$ and
$2$ (values $1+3=4$), which beats robbing house $1$ alone or house $3$
alone or houses $1,3$ together ($2+1=3$).

\begin{complexitybox}[Approach 3 Complexity]
\textbf{Time.} $O(n)$.
\textbf{Space.} $O(n)$ table, no recursion stack.
\end{complexitybox}

\subsection*{Approach 4: Space Optimization}

\begin{lstlisting}
#include <bits/stdc++.h>
using namespace std;

int robSpaceOptimized(vector<int>& nums) {
    int n = nums.size();
    int prev2 = 0, prev1 = nums[0];

    for (int i = 1; i < n; i++) {
        int skip = prev1;
        int take = prev2 + nums[i];
        int current = max(skip, take);
        prev2 = prev1;
        prev1 = current;
    }
    return prev1;
}

int main() {
    vector<int> nums = {1, 2, 3, 1};
    cout << robSpaceOptimized(nums) << endl; // 4
    return 0;
}
\end{lstlisting}

\begin{complexitybox}[Approach 4 Complexity]
\textbf{Time.} $O(n)$.
\textbf{Space.} $O(1)$.
\end{complexitybox}

\begin{takeawaybox}
``Maximum sum subject to a no-two-adjacent constraint'' is solved by an
include-or-exclude recurrence at every position, maximized rather than
summed -- the DP analogue of the pick/not-pick backtracking tree from
the Recursion chapter's Power Set section, but resolved in $O(n)$ instead of
$O(2^n)$ because we only need the \emph{best} outcome at each position,
not every possible outcome. Section~\ref{sec:dp-house-robber-2} adapts
this exact recurrence to a circular arrangement of houses.
\end{takeawaybox}

\section{House Robber II (Houses in a Circle)}
\label{sec:dp-house-robber-2}

\problemstatement{Same as Section~\ref{sec:dp-house-robber}, except the
houses are arranged in a \textbf{circle} -- house $0$ and house $n-1$ are
now also considered adjacent, so robbing both of them together is
forbidden, in addition to the usual no-two-consecutive-houses rule.}

\begin{patternbox}
\emph{``Same problem, but the arrangement is circular instead of
linear''} is a keyword for a \textbf{reduction} technique: rather than
deriving a new recurrence that somehow accounts for a wraparound
adjacency, we instead run the \emph{already-solved} linear House Robber
(Section~\ref{sec:dp-house-robber}) twice, on two carefully chosen
linear sub-ranges, and combine their results.
\end{patternbox}

\subsection*{Understanding the problem carefully}

The only new constraint, compared to the linear version, is that house
$0$ and house $n-1$ cannot \emph{both} be robbed. Equivalently: in any
valid circular solution, \textbf{at least one} of house $0$ or house
$n-1$ is never robbed. This means every valid circular solution is also
a valid \emph{linear} solution either to the sub-array excluding house
$0$ (i.e.\ houses $1$ through $n-1$), or to the sub-array excluding
house $n-1$ (i.e.\ houses $0$ through $n-2$) -- because whichever of the
two ``forbidden together'' houses is excluded from the solution can
simply be excluded from consideration entirely.

\begin{ideabox}[Reduce to Two Linear Sub-problems, Take the Better One]
Run Section~\ref{sec:dp-house-robber}'s linear House Robber on the
sub-array $\textit{nums}[0..n-2]$ (excluding the last house), and again
on $\textit{nums}[1..n-1]$ (excluding the first house). The answer to the
circular version is the \textbf{maximum} of these two results.
\end{ideabox}

\subsection*{Why considering exactly these two sub-ranges is sufficient}

Every valid circular solution falls into (at least) one of two
categories: those that do not rob house $n-1$ (entirely captured by the
linear sub-problem on houses $0$ through $n-2$, since within that range
the usual linear adjacency rule is the \emph{only} remaining constraint),
and those that do not rob house $0$ (captured by the linear sub-problem
on houses $1$ through $n-1$). A solution that robs \emph{neither} house
$0$ nor house $n-1$ is captured by both sub-problems redundantly, which
costs nothing -- we just take whichever of the two computed maximums is
larger, and that larger value is guaranteed to be a valid circular
answer that misses no better possibility, since any valid circular
solution at all must belong to at least one of these two categories.

\subsection*{Approach 1: Recursion (via Two Calls to the Linear Helper)}

\begin{lstlisting}
#include <bits/stdc++.h>
using namespace std;

int robRecursive(int i, vector<int>& nums) {
    if (i == 0) return nums[0];
    if (i < 0) return 0;

    int skip = robRecursive(i - 1, nums);
    int take = robRecursive(i - 2, nums) + nums[i];
    return max(skip, take);
}

int robLinear(vector<int> sub) {
    if (sub.empty()) return 0;
    return robRecursive((int)sub.size() - 1, sub);
}

int robCircular(vector<int>& nums) {
    int n = nums.size();
    if (n == 1) return nums[0];

    vector<int> excludeLast(nums.begin(), nums.end() - 1);
    vector<int> excludeFirst(nums.begin() + 1, nums.end());

    return max(robLinear(excludeLast), robLinear(excludeFirst));
}

int main() {
    vector<int> nums = {2, 3, 2};
    cout << robCircular(nums) << endl; // 3
    return 0;
}
\end{lstlisting}

\begin{complexitybox}[Approach 1 Complexity]
\textbf{Time.} $O(2^n)$ for each of the two linear calls, so $O(2^n)$
overall (the constant factor of $2$ calls does not change the
asymptotic class).
\textbf{Space.} $O(n)$ recursion stack.
\end{complexitybox}

\subsection*{Approaches 2--4: Memoization, Tabulation, and Space Optimization}

Each of these three later stages is obtained by substituting
Section~\ref{sec:dp-house-robber}'s corresponding memoized, tabulated, or
space-optimized linear House Robber implementation in place of
$\textit{robLinear}$ above -- no new derivation is required, since the
reduction itself (call the linear solver twice, take the max) is
independent of which of the four stages that linear solver internally
uses.

\begin{lstlisting}
#include <bits/stdc++.h>
using namespace std;

int robLinearSpaceOptimized(vector<int>& nums) {
    if (nums.empty()) return 0;
    int n = nums.size();
    int prev2 = 0, prev1 = nums[0];

    for (int i = 1; i < n; i++) {
        int skip = prev1;
        int take = prev2 + nums[i];
        int current = max(skip, take);
        prev2 = prev1;
        prev1 = current;
    }
    return prev1;
}

int robCircularSpaceOptimized(vector<int>& nums) {
    int n = nums.size();
    if (n == 1) return nums[0];

    vector<int> excludeLast(nums.begin(), nums.end() - 1);
    vector<int> excludeFirst(nums.begin() + 1, nums.end());

    return max(robLinearSpaceOptimized(excludeLast),
               robLinearSpaceOptimized(excludeFirst));
}

int main() {
    vector<int> nums = {2, 3, 2};
    cout << robCircularSpaceOptimized(nums) << endl; // 3
    return 0;
}
\end{lstlisting}

\subsection*{Dry run}

\dryrun{Take $\textit{nums} = [2,3,2]$.}

\begin{itemize}
  \item Exclude last (sub-array $[2,3]$): linear House Robber gives
        $\max(2,3)=3$.
  \item Exclude first (sub-array $[3,2]$): linear House Robber gives
        $\max(3,2)=3$.
  \item Overall answer: $\max(3,3)=3$.
\end{itemize}

Checking by hand: robbing houses $0$ and $2$ together is forbidden
(circular adjacency), so the best options are house $1$ alone ($3$), or
house $0$ alone ($2$), or house $2$ alone ($2$) -- confirming $3$ is
correct.

\begin{complexitybox}[Final (Space-Optimized) Complexity]
\textbf{Time.} $O(n)$ -- two linear passes, each $O(n)$.
\textbf{Space.} $O(1)$ extra (beyond the two $O(n)$ sliced sub-arrays,
which could themselves be avoided by passing index ranges instead of
copying, reducing auxiliary space to true $O(1)$).
\end{complexitybox}

\begin{takeawaybox}
When a new constraint (here, circularity) only affects the boundary
between two elements, check whether the problem can be split into a
small, fixed number of cases that each remove the offending boundary
entirely -- reducing the new problem to multiple calls of an
already-solved simpler one, rather than re-deriving a new recurrence
from scratch. This reduction technique is often faster to get correct
than attempting to encode the wraparound constraint directly into a
single, more complex recurrence.
\end{takeawaybox}

